\section{Introduction}
\label{sec:intro}

Retrieval-augmented generation (RAG) decomposes complex questions into
retrieval and reasoning steps, and a growing line of work makes those steps
increasingly explicit: logical query trees \cite{planrag}, executable programs
\cite{pyrag}, typed slots \cite{slotrag}, and learned adaptive control
\cite{dynakrag}. This explicitness is a prerequisite for treating multi-hop
retrieval as an \emph{execution} problem rather than a prompt-engineering
problem. But it also raises a systems question that most RAG papers do not
state precisely: given a finite resource budget, \emph{how should a system
compile a question into dependent evidence requirements, choose physical
implementations for them, allocate budget across them, and acquire the evidence
while observing the outcomes at runtime}?

This paper answers that question with a concrete evidence-centric execution
model. We make three connected moves. First, we treat evidence \emph{requirements}
--- conditions that must be satisfied, with a type, an importance, a set of
variables to bind, and dependencies over other requirements --- as first-class
objects, and give them a typed algebra with monotone state-transition
semantics (\S\ref{sec:algebra}). Second, we separate the logical evidence
program from its physical implementations and introduce a deterministic
importance-weighted allocator that assigns execution orders, retrieval
implementations, and budget allocations to typed requirements under a matched
budget (\S\ref{sec:optimizer}). Third, we execute plans deterministically with
hard budget enforcement and full provenance, so that every answer is traceable
to the evidence and physical actions that produced it (\S\ref{sec:execution}).

The design is governed throughout by an honesty discipline that we apply to
ourselves. A property estimator is claimed only where it is construction-derived
and exact; on the well-defined chains that dominate our workloads, requirement
importance equals $\tau = 2d-1$ by derivation, and we do not claim a learned
estimator adds value there. Runtime re-optimization over alternative physical
plans is explicitly future work under a branching execution model, because on
the chain-only topology this system runs today it would be vacuous. And in the
evaluation, a win is claimed only when it is statistically supported under the
matched-budget protocol; point-estimate differences are reported as such.

Our main experimental finding is that the requirement-aware optimizer, while
modest in its absolute accuracy footprint, is a \emph{cost}-oriented win on
the subdomain where it has degrees of freedom, with directionally positive
EM. On the full matched development sample it does not spend more on average
than the frozen static baseline (pooled $-0.16$ calls/question; HotpotQA $-0.19$, 2Wiki $-0.24$,
MuSiQue $-0.09$), and EM differences are directionally positive on all three
datasets ($\Delta$EM $+6.3/+5.9/+9.1$pt). Neither effect is statistically
supported at this sample size: cost reductions have bootstrap
$p \in [0.119,0.376]$ and EM differences have McNemar $p \geq 0.625$ (pooled
EM $+7.3$pt on 55 matched questions). On the $\ge 3$-slot subdomain, where the
optimizer has real allocation freedom, the cost saving is larger and
statistically supported in the cleanest comparison: pooled $3.45$ calls/question
versus $4.73$ for the cost-only arm ($\Delta=-1.27$, bootstrap $p<0.001$;
vs.\ static $-1.09$, $p=0.036$; a deterministic structural finding on
$\ge 3$-slot chains), at pooled EM that is equal across arms
($54.5\%/45.5\%/54.5\%$). We make no quality claim on EM. The honest boundary is
that the $1$--$2$-slot chains dominating most QA splits leave the optimizer
little allocation freedom, so the cost saving is concentrated on the
$\ge 3$-slot subdomain.

\paragraph{Contributions.}
\begin{itemize}
\item An evidence-centric execution model for multi-hop RAG, in which a
question is compiled into typed, dependent evidence requirements with a formal
monotone state-transition algebra (\S\ref{sec:algebra}).
\item A requirement-aware physical-plan allocator that assigns legal orders,
physical implementations, and budget allocations to typed requirements under a
matched budget, derived from the chain-law importance
$\tau = 2d-1$ on well-defined chains (\S\ref{sec:optimizer}).
\item A deterministic adaptive executor with hard budget enforcement and
end-to-end provenance, enabling per-question trace audits and operator-level
failure attribution (\S\ref{sec:execution}).
\item An honest evaluation under a matched-budget protocol: a uniform cost
reduction without any quality regression on the full development sample
(pooled $-0.16$ calls/question, directionally positive EM $+7.3$pt; neither
statistically supported at this sample size, reported as such), and a
statistically supported $\ge 3$-slot subdomain saving of $1.27$ calls/question
over the cost-only arm (bootstrap $p<0.001$) at equal pooled EM---defining the
boundary where the optimizer has little allocation freedom
(\S\ref{sec:results}).
\end{itemize}
